% This file should be rename with the speaker's last name. (e.g : smith)

% Write the author names in the order you wish them to appear
% and underline the one who will give the talk.

% Only the speaker's email address is required.

\begin{abstract_online}{Teaching Decision Analysis using a Computer Algebra System}{%
        \underline{Karsten Schmidt}$^{1}$}{[kschmidt@hs-sm.de]
        }{%
        $^1$ Schmalkalden University of Applied Sciences, Germany \newline}
    
In the Faculty of Business and Economics at Schmalkalden University, the Decision Analysis course in the bachelor program is routinely taught in a traditional classroom setting (blackboard, overhead projector, and pocket calculators). This course is actually one half of the subject “Mathematics II”, the other half is Matrix Algebra, which has been taught in the PC lab for many years (one or two students in front of a PC, instructor’s PC connected to a projector). As the teacher of the Decision Analysis course is currently on maternity leave, I took over teaching of this course for two years from her.

I was curious if topics from the Matrix Algebra portion of “Mathematics II” were useful in the Decision Analysis portion. Particularly as in decision analysis a large number of matrices (sometimes called tables) is used, for example payoff matrices, results matrices, harm matrices, opportunity costs matrices. However, the answer is No.

Nevertheless, having a Computer Algebra System (CAS) readily available is not only useful for matrix operations, but also for finding the perfect alternative, or action, in a decision problem, using other mathematical methods. The usefulness of a CAS in decision analysis will be demonstrated in several examples from different areas, e.g. decisions under certainty, decisions under uncertainty, and decisions under risk.

As the students learn to work with the CAS in the matrix algebra portion anyway, using it (together with a spreadsheet program) also in the decision analysis portion comes without a steep learning curve. Note that students can install the CAS legally on their private PCs as long as they are enrolled in our faculty, and have access to it during the final exam in the PC lab (then, naturally, only one student per PC).



\end{abstract_online}
    